\subsection{Sugestões para trabalhos futuros}
\label{sec:cap5_trabalhos_futuros}

Os resultados obtidos fornecem base consistente para o estudo de operações de \gls{rvdb} com manipulador robótico em uma espaçonave robótica, mas também indicam possibilidades de extensão do trabalho.
A seguir, apresentam-se sugestões para estudos futuros.

\subsubsection{Extensões do modelo geométrico e dinâmico}

\begin{enumerate}[label=(\alph*)]

    \item \textbf{Modelagem de contato e acoplamento}

    Introduzir modelos de contato entre ferramenta e porta, incluindo rigidez, amortecimento e elementos geométricos, como cones de engate ou guias mecânicas.
    Com isso, seria possível:
    \begin{itemize}
        \item estudar a fase de \emph{soft capture} com controle de força ou impedância no manipulador;
        \item analisar a transferência de impulso e as oscilações induzidas na base durante o acoplamento;
        \item definir tolerâncias posicionais e angulares mais realistas para a interface de engate.
    \end{itemize}

    \item \textbf{Dinâmica orbital com perturbações ambientais e achatamento da Terra}

    Estender o modelo orbital com a inclusão de perturbações, como o termo de achatamento $J_2$, arrasto residual em órbitas baixas e incertezas paramétricas.
    Essa extensão permitiria:
    \begin{itemize}
        \item quantificar a robustez do controlador translacional diante de perturbações sistemáticas;
        \item avaliar desvios acumulados em missões de maior duração;
        \item investigar a necessidade de leis de controle mais sofisticadas para manter as janelas de posição e velocidade.
    \end{itemize}

\end{enumerate}

\subsubsection{Estratégias de controle}

\begin{enumerate}[label=(\alph*),resume]

    \item \textbf{Controle adaptativo ou com ganho variável}

    Investigar controladores com ganhos variando em função da fase da aproximação (por exemplo, distância ao alvo, velocidade relativa, energia cinética ou outros indicadores).
    Essa abordagem permitiria:
    \begin{itemize}
        \item reduzir tempos de estabilização em trechos específicos da missão;
        \item ajustar automaticamente os ganhos de controle conforme o regime dinâmico.
    \end{itemize}

    \item \textbf{Análise formal de estabilidade e robustez}

    Desenvolver análises formais de estabilidade e robustez para o sistema acoplado base--manipulador, por exemplo:
    \begin{itemize}
        \item construção de funções de Lyapunov apropriadas para o sistema completo;
        \item uso de critérios no domínio da frequência para avaliar margens de estabilidade;
        \item derivação de condições de estabilidade diante de variações de massa, inércia, atrasos e perturbações externas.
    \end{itemize}
    Esses estudos forneceriam garantias mais rigorosas para o uso da arquitetura de controle em cenários operacionais mais exigentes.

\end{enumerate}
